\subsection{积分中取极限}

命题 3. --- 设 B 是 \(L_{F}^{1}\) 上的一个滤子基。假定存在紧集 \(K \subset X\) ，使对每个集 \(M \in B\) ，所有函数 \(f \in M\) 的支集都含于 K。在这些条件下，若 B 在 X 上一致收敛于 \(f_{0}\) ，则函数 \(f_{0}\) 是可积的，且

\[
\int \mathbf {f} _ {0} d \mu = \lim _ {\mathfrak {B}} \int \mathbf {f} d \mu .\tag{9}
\]

因为，\(\mathfrak{B}\) 平均收敛于 \(\mathbf{f}_0\) （§3, No.~3, 命题 4）。

命题 4. --- 设 \((f_n)\) 是可积数值函数的一个递增（相应地，递减）序列。为使序列的上（相应地，下）包络 \(f\) 可积，必须且只须 \(\sup_{n} \int f_n d|\mu| < +\infty\) （相应地 \(\inf_{n} \int f_n d|\mu| > -\infty\) ）；此时

\[
\int f d \mu = \lim _ {n \to \infty} \int f _ {n} d \mu .\tag{10}
\]

我们只考虑递增序列。序列 \(g_{n} = f_{n} + f_{1}^{-}\) 是递增的，由可积函数 \(\geqslant 0\) 组成；由于它的上包络是 \(g = f + f_{1}^{-}\) ，命题由 §3, No.~6 的定理 5 得出。

定理 2. --- 设 A 是一个指标集，由具有可数基的滤子 F 过滤。设 \((\mathbf{f}_{\alpha})_{\alpha\in\mathrm{A}}\) 是可积函数的一个族，它关于滤子 F 几乎处处逐点收敛于一个函数 f；若存在数值函数 \(g\geqslant0\) 使 \(\int^{*}gd|\mu|<+\infty\) 且 \(|\mathbf{f}_{\alpha}(x)|\leqslant g(x)\) 在 X 中几乎处处成立对每个 \(\alpha\in A\) ，则函数 f 是可积的，且

\[
\int \mathbf {f} d \mu = \lim _ {\mathfrak {F}} \int \mathbf {f} _ {\alpha} d \mu .\tag{11}
\]

定理由 Lebesgue 定理（§3, No.~7, 定理 6 的推论）得出，因为在陈述的条件下，\(\mathbf{f}_{\alpha}\) 平均收敛于 \(\mathbf{f}\) （关于 \(\mathfrak{F}\) ）。

推论 1. --- 设 \(\Omega\) 是一个拓扑空间，\(t_0\) 是 \(\Omega\) 的一个具有可数基本邻域系的点，\(\mathbf{f}\) 是 \(\mathbf{X} \times \Omega\) 到 \(\mathbf{F}\) 中的一个映射，具有下列性质：

\begin{enumerate}
\def\labelenumi{\alph{enumi})}
\item
  对每个 \(t \in \Omega\) ，函数 \(x \mapsto \mathbf{f}(x, t)\) 是可积的；
\item
  对每个 \(x \in \mathrm{X}\) ，函数 \(t \mapsto \mathbf{f}(x, t)\) 在 \(t_0\) 处连续；
\item
  存在 \(t_0\) 的一个邻域 U 与定义在 X 上的数值函数 \(g \geqslant 0\) ，使 \(\int^{*} g d|\mu| < +\infty\) 且 \(|\mathbf{f}(x,t)| \leqslant g(x)\) 对所有 \(x \in X\) 与 \(t \in U\) 成立。
\end{enumerate}

在这些条件下，映射 \(t \mapsto \int \mathbf{f}(x,t) d\mu(x)\) （\(\Omega\) 到 \(\mathbf{F}\) 中的）在点 \(t_0\) 处连续。

推论 2. --- 设 \((\mathbf{f}_{n})\) 是可积函数的一个序列，使通项为 \(\mathbf{f}_{n}(x)\) 的级数几乎处处收敛；若存在函数 \(g \geqslant 0\) 使 \(\int^{*} g d|\mu| < +\infty\) 且对每个整数 n 有 \(\left|\sum_{k=1}^{n} \mathbf{f}_{k}(x)\right| \leqslant g(x)\) 几乎处处成立，则和 \(\mathbf{f}(x)\) （通项为 \(\mathbf{f}_{n}(x)\) 的级数的，几乎处处有定义的）是可积的，且

\[
\int \mathbf {f} d \mu = \sum_ {n = 1} ^ {\infty} \int \mathbf {f} _ {n} d \mu\tag{12}
\]

（«级数的逐项积分»）。

